\documentclass[11pt,a4paper]{article}

\usepackage[utf8]{inputenc}
\usepackage[T1]{fontenc}
\usepackage[spanish,es-noquoting]{babel}
\usepackage{amsmath,amssymb}
\usepackage{graphicx}
\usepackage{siunitx}
\usepackage{booktabs}
\usepackage[hidelinks]{hyperref}
\usepackage[margin=2.5cm]{geometry}

\title{Notes}
\author{Andrés}
\date{\today}

\begin{document}
\newpage
\section{Notes}
I used the Griffiths particle physics book for theory of these notes.

\section{Weak forces}
To start with the weak interaction we need to know, all the particles can be used in these interactions.\\
The two kinds of weak interactions are:
\begin{itemize}
    \item Neutral
    \item Charged
\end{itemize}
The interaction of this force occurs when a particle changes to another type. An example of this is when a up quark wants to change to a down one, this interaction involves the use of a boson, in this case the $W^{+}$ boson .\\
There are three types of bosons for these interactions, two that carry the electric charges and one that electric charge is neutral
\begin{itemize}
    \item $W^+$ Electric charge: +1
    \item $W^-$ Electric charge: -1
    \item $Z$ Electric charge: 0
\end{itemize}
Now I can explain the properties and the information of each boson independently
\subsection{Boson Z}
As I explain in the past section this particle has a neutral charge, that implies the interactions where is used implies there are some "liberty", because the transofrmations can make are without the implication of the electric charge.\\
The Z boson can also be in some way similar to the use of the electromagentic. Imagine we have to electrons that are repell, we know this repell is caused by the photon boson, but if we make the arrange for the Z boson:\\
\begin{equation}
    e^- + e^- \rightarrow e^-+e^-
\end{equation}
Because the Z boson doesn't have an electric charge in this interaction can also change the momentum of the electrons. The explanation of this phenomenon are in the most cases are explain by the electric force because the Z boson has mass; that implies that the propagation have a mass coefficent. For high energies the interaction made by the Z boson are stronger than the photon.\\
Taking the fact the Z boson have mass we can say the exact amount: 91.2 $GeV/c^2$
\end{document}
